% =============================================================================
%  main.tex -- venue-agnostic paper skeleton (article class)
%
%  What this is
%    A self-contained skeleton that compiles with a stock TeX Live install.
%    It is NOT a venue template: it has no venue-specific layout, page limit or
%    reference style. For IEEE / ACL / NeurIPS, read README.md and the stubs in
%    ieee/ and acl/.
%
%  How to use
%    1. Replace every \fillme{...} box with real prose. Nothing else in this
%       file is allowed to be vague at submission time.
%    2. Resolve every "% TODO(...)" comment (`grep -n TODO main.tex`).
%    3. Never type a number by hand. Every number must come from
%       04_experiments/RESULTS-INDEX.json and carry a "% num:<id>" comment on
%       the line immediately above it (tools/CONTRACTS.md section 3).
%    4. Never add a citation that is not in 02_literature/verified_citations.csv
%       (RULES.md #1). refs.bib ships with zero active entries on purpose.
%    5. Tables and figures are generated, not typed: run
%          python tables/make_tables.py  ...
%          python figures/make_figures.py ...
%       The \input blocks below pick the generated files up automatically and
%       fall back to a visible placeholder when they are missing.
%
%  Compile
%    ./build.sh          (needs latexmk + TeX Live; see BUILD.md)
%    This file has NOT been compiled in the repo's build environment (no LaTeX
%    installed there). See BUILD.md, "Verified status".
% =============================================================================

\documentclass[11pt,a4paper]{article}

\usepackage[T1]{fontenc}
\usepackage[utf8]{inputenc}
\usepackage[margin=1in]{geometry}
\usepackage{amsmath,amssymb}
\usepackage{graphicx}
\usepackage{booktabs}
\usepackage[font=small,labelfont=bf]{caption}
\usepackage[numbers,sort&compress]{natbib}
\usepackage{xcolor}
\usepackage{url}
\usepackage[hidelinks]{hyperref}

% -----------------------------------------------------------------------------
% Macros. Define the method name once; rename it here only.
% TODO(name): replace with the method's real name. Do not name a method after
% a dataset or benchmark you did not build.
% -----------------------------------------------------------------------------
\newcommand{\method}{TODO-Method-Name}

% Visible placeholder. Every \fillme{...} must be gone before submission;
% tools/check_numbers.py and the repo linter grep for it.
\newcommand{\fillme}[1]{%
  \par\noindent\fbox{\parbox{0.96\linewidth}{\small\textbf{TODO} --- #1}}\par}

% Fenced code / command in running text.
\newcommand{\code}[1]{\texttt{#1}}

% -----------------------------------------------------------------------------
% Title block
% -----------------------------------------------------------------------------
\title{%
  TODO(title): a specific, falsifiable title.\\
  \large TODO(subtitle, optional): the claim in one clause, no adjectives.%
}

\author{%
  TODO(Author One)\thanks{Corresponding author: TODO(email).} \and
  TODO(Author Two)\thanks{%
    \textbf{Use of AI tools.} TODO(disclosure): name each tool and version, the
    sections it touched, and what it was used for (e.g.\ language editing of
    Section~\ref{sec:related}, or an outline the author rewrote). This
    disclosure is written to satisfy the target venue's policy on AI-assisted
    writing; TODO(venue): name the venue, quote the policy URL and add the
    date it was retrieved (starting points: 00\_foundations/ETHICS-AND-AI-USE.md).
    TODO(verify at drafting time): the following three sentences are the target
    state, not a verified fact about this paper --- do not let them stand unless
    each is true on the day of submission, and delete any that is not.
    (i) no AI tool produced data, results, figures or references;
    (ii) every citation was resolved and recorded in
    \code{02\_literature/verified\_citations.csv};
    (iii) the tools are not authors because they cannot take responsibility for
    the work. See Section~\ref{sec:ethics}.}%
}

\date{TODO(date): submission date, or leave empty for \texttt{\textbackslash today}}

\begin{document}
\maketitle

% -----------------------------------------------------------------------------
\begin{abstract}
% TODO(abstract): 150-250 words, written last, containing exactly: the problem,
% what was done, the headline number (marked with % num:<id>), and the main
% limitation. No citations, no undefined acronyms, no "state of the art".
\fillme{Abstract: problem (1 sentence) / method (1 sentence) / headline result
with its index id / limitation (1 sentence).}
\end{abstract}

% -----------------------------------------------------------------------------
\section{Introduction}
\label{sec:intro}

% TODO(intro-p1): the problem, in plain language, for a reader outside the
% subfield. One concrete example beats three general sentences.
% TODO(intro-p2): what is missing in existing work -- this sets up
% Section~\ref{sec:related}. Cite only verified keys.
% TODO(intro-p3): what this paper does, one sentence per contribution, each
% contribution mapped to a section or a table.
% TODO(intro-p4): the headline result, with its "% num:<id>" marker, and the
% single most important caveat (do not bury it in Section~\ref{sec:limitations}).

\fillme{Introduction: problem / gap / contributions / headline result.}

% -----------------------------------------------------------------------------
\section{Related Work}
\label{sec:related}

% TODO(related): organise by *idea*, not by paper. For each cluster: what the
% cluster assumes, where it breaks, and how this work differs. Every sentence
% that names prior work must carry a \citep{} whose key exists in
% 02_literature/verified_citations.csv. If you cannot verify a citation, drop
% the sentence (RULES.md #1).
% TODO(related-delta): one explicit paragraph "how this differs", because a
% reviewer will look for it first.

\fillme{Related Work: 2-4 thematic clusters, then an explicit delta paragraph.}

% -----------------------------------------------------------------------------
\section{Method}
\label{sec:method}

% TODO(method-setup): notation table or a short paragraph fixing symbols.
% TODO(method-core): the method itself, in the order a reader must understand
% it. Every design choice needs a reason or a pointer to an ablation
% (Table~\ref{tab:ablation}).
% TODO(method-complexity): cost, in time and memory, measured -- not estimated.

\fillme{Method: setup and notation, the method, and its cost.}

\begin{equation}
\label{eq:objective}
  \mathcal{L}(\theta) = \underbrace{\frac{1}{N}\sum_{i=1}^{N}\ell\big(f_\theta(x_i), y_i\big)}_{\text{data term}}
  + \lambda\,\underbrace{\Omega(\theta)}_{\text{regulariser}}
\end{equation}

% TODO(eq:objective): replace Equation~\ref{eq:objective} with the real
% objective, or delete it. Keep the label if other text references it.

% -----------------------------------------------------------------------------
\section{Experiments}
\label{sec:exp}

% TODO(exp-data): datasets, splits, sizes, and the licence/access route for each
% (public data only -- RULES.md #6).
% TODO(exp-baselines): each baseline and why it is the right comparison; state
% whether you re-ran it or copied the number, and if copied, from where.
% TODO(exp-protocol): seeds, number of runs, hardware, wall-clock, and the exact
% command. This paragraph must match 04_experiments/PROTOCOL.md.
% TODO(exp-metrics): what is measured and in which direction is better.

\fillme{Experiments: data / baselines / protocol (seeds, runs, command) / metrics.}

\subsection{Implementation details}
\label{sec:impl}

% TODO(impl): framework versions, hyperparameters (Appendix~\ref{app:hparams}),
% and the commit hash the numbers came from.

\fillme{Implementation details: versions, hyperparameters, commit hash.}

% -----------------------------------------------------------------------------
\section{Results}
\label{sec:results}

% TODO(results-p1): answer the research question in one paragraph before any
% table. Then point at Table~\ref{tab:main-results} and
% Figure~\ref{fig:main-accuracy}. Numbers come from the generated files; do not
% retype them.
% TODO(results-p2): the ablation (Table~\ref{tab:ablation}) -- what each
% component contributes, and what did NOT help.
% TODO(results-negative): negative and null results belong here, stated plainly.

\fillme{Results: the answer first, then the evidence, then what did not work.}

% --- Table 1: generated by tables/make_tables.py from a summary_*.json -------
% The caption and label live here so that the generated file stays a pure
% tabular body and labels cannot collide. Paths are relative to this directory.
\begin{table}[t]
  \centering
  \caption{%
    TODO(caption): what is compared, on what data, over how many seeds, and what
    the interval means. Example shape: ``Test accuracy on TODO(dataset);
    mean $\pm$ standard deviation over $n$ seeds; 95\% CI is a bootstrap
    interval.'' Every number in this table is generated from
    \code{04\_experiments/results/summary\_TODO.json}.}
  \label{tab:main-results}
  \IfFileExists{tables/generated/main-results.tex}{%
    \input{tables/generated/main-results.tex}%
  }{%
    \fbox{\parbox{0.9\linewidth}{\small No generated table yet. Run
      \code{python tables/make\_tables.py --summary
      04\_experiments/results/summary\_TODO.json --out-dir tables/generated
      -{}-name main-results}.}}%
  }
\end{table}

% --- Figure 1: generated by figures/make_figures.py --------------------------
\begin{figure}[t]
  \centering
  \IfFileExists{figures/generated/main-accuracy.tex}{%
    \input{figures/generated/main-accuracy.tex}%
  }{%
    \fbox{\parbox[c][4cm][c]{0.8\linewidth}{\centering\small No generated
      figure yet. Run \code{python figures/make\_figures.py --index
      04\_experiments/RESULTS-INDEX.json --figure main-accuracy
      -{}-out-dir figures/generated}.}}%
  }
  \caption{%
    TODO(caption): what the axes are, what the bars/points are, and the one
    thing the reader should notice. State the unit and the uncertainty.}
  \label{fig:main-accuracy}
\end{figure}

% --- Table 2: hand-written example of the booktabs style ---------------------
% Keep this table only if you have an ablation; otherwise delete it and its
% reference in the text. Placeholder cells say TODO on purpose -- no invented
% numbers are allowed anywhere in this repo.
\begin{table}[t]
  \centering
  \caption{%
    TODO(caption): ablation -- which component is removed in each row and what
    that costs. Same data, same seeds, same protocol as
    Table~\ref{tab:main-results}.}
  \label{tab:ablation}
  \begin{tabular}{lcc}
    \toprule
    Variant & TODO(metric) & $\Delta$ vs.\ full \\
    \midrule
    \method{} (full)          & TODO(value) & --- \\
    \hspace{1em}-- without TODO(part A) & TODO(value) & TODO(delta) \\
    \hspace{1em}-- without TODO(part B) & TODO(value) & TODO(delta) \\
    \bottomrule
  \end{tabular}
\end{table}

% -----------------------------------------------------------------------------
\section{Limitations}
\label{sec:limitations}

% TODO(limitations): the honest list. For each: what it is, who it affects, and
% what would fix it. Cover at least: data coverage, compute budget, number of
% seeds and what that means for the confidence intervals, and any claim in the
% abstract that the evidence only partly supports. Reviewers read this section
% first; an empty one is a rejection reason.

\fillme{Limitations: data, compute, statistics, scope. Name the claim you cannot
fully support.}

% -----------------------------------------------------------------------------
\section{Ethics Statement}
\label{sec:ethics}

% TODO(ethics): (a) human subjects / personal data -- say ``none'' if none;
% (b) licences of every dataset and model used; (c) foreseeable misuse;
% (d) the AI-assistance disclosure, which must match the footnote on the title
% page exactly, and must follow the target venue's policy rather than a generic
% one (IEEE and ACL put it in the Acknowledgements; ACM asks in the methods and
% only for research use; NeurIPS asks in the checklist). Sources:
% 00_foundations/ETHICS-AND-AI-USE.md.

\fillme{Ethics: data and licences, misuse, and the AI-assistance statement
(cross-reference the footnote on page~1; keep the two texts identical).}

% -----------------------------------------------------------------------------
\section{Reproducibility Statement}
\label{sec:repro}

% TODO(repro): where the code, configs and run logs live; the exact commands;
% the seeds; what is released and what cannot be (and why). Give one command
% that reproduces Table~\ref{tab:main-results}, and state the tolerance at
% which a rerun counts as reproducing the number.

\fillme{Reproducibility: repo path, command, seeds, released artefacts, rerun
tolerance.}

% -----------------------------------------------------------------------------
\section{Conclusion}
\label{sec:conclusion}

% TODO(conclusion): what is now known that was not known before, in two
% sentences, calibrated to the evidence (no ``state of the art''; RULES.md #5).
% Then the single most useful next experiment.

\fillme{Conclusion: what is now known, and the next experiment worth running.}

% -----------------------------------------------------------------------------
% References. refs.bib ships with ZERO active entries: entries may only be
% added from 02_literature/verified_citations.csv. To populate it:
%   python tools/citations_to_bib.py --csv 02_literature/verified_citations.csv \
%          --out refs.generated.bib
% then either append refs.generated.bib to refs.bib or change the line below to
%   \bibliography{refs,refs.generated}
% An empty bibliography produces "no \citation" warnings -- that is expected
% for an empty skeleton and not an error.
\bibliographystyle{plainnat}
\bibliography{refs}

% -----------------------------------------------------------------------------
\appendix
\section{Hyperparameters}
\label{app:hparams}

% TODO(appendix): every hyperparameter, its search range, and the selected
% value. Copy from the config JSON, do not retype from memory.

\fillme{Appendix: full hyperparameter table, search ranges, selected values.}

\end{document}
